%%%%%%%%%%%%%%%%%%%%%%%%%%%%%%%%%%%%%%%%%%%%%%%%%%%%%%%%%%%%
\section{引言}
\label{sec:introduction}

企业正苦于将智能体 AI（agentic AI）系统部署到生产环境。这些系统承诺能自动化复杂工作流——管理金融交易、病历记录和关键基础设施——但它们引入的安全挑战是现有防御手段无法应对的。例如，在发布 OpenAI Atlas 浏览器后数小时内，研究人员就演示了网页内容中嵌入的恶意指令可触发助手泄露凭证~\cite{atlas-cve}。OpenAI 的 CISO 承认"提示注入（prompt injection）仍然是一个未解决的问题"~\cite{openai-ciso-atlas}。

挑战远不止提示注入本身。LLM 无法区分指令与数据。非确定性执行路径无法预测或静态分析。多轮交互允许渐进式上下文操控，其中每条消息看似无害。现有方法基于模式查找或概率判断——两者都需要枚举攻击，使防御者困于未知之未知。

本方案通过系统化方法消除未知之未知：不是枚举攻击，而是确定性界定系统。核心洞察是：智能体系统可抽象为简单、拜占庭式的分布式系统——多个实体在明确界定的边界上交互。安全问题遂归结为保护边界穿越。概念上，本方案基于在每个边界穿越处满足两个属性：理解\textbf{意图（intent）}以确保操作满足组织策略，强制执行\textbf{完整性（integrity）}以确保操作真实且未被篡改。两者缺一不可，单独任意一个都不充分。
此设计将智能体工作流约简为确定性分布式系统，其中边界由策略守卫，穿越由密码学保护。攻破系统需要攻破密码学，而非精心构造提示。多类攻击——如身份伪造、会话重放、策略替换等——在设计层面（by-design）即被消除：它们必须攻破密码学原语才能得逞；其他攻击——如未授权数据访问、权限提升、凭证泄露等——则在运行时由策略阻断。
架构上，\textbf{策略层面强制执行（by-policy enforcement）}与\textbf{设计层面消除（by-design elimination）}共同提供互补、可组合的纵深防御。

我们研究了 100+ 个智能体应用，识别出关键挑战：\textbf{(1) 四个并发攻击面}——工具、数据、提示、上下文——攻击可跨面组合，达成任意单一表面无法实现的目标。\textbf{(2) 异构框架}——智能体构建使用多种协议（MCP、A2A）、LLM 接口（Claude、OpenAI）、编排器（LangChain、CrewAI、AutoGen）。异构性是生态系统的持久特征。\textbf{(3) 组合式安全缺口}——逐框架检查会遗漏跨越组合的违规。\textbf{(4) 开发者负担}——要求在异构框架的每一次交互中嵌入安全逻辑是不现实的。

综合来看，攻击面的增长速度超过工具的应对速度，凸显出一个关键缺口：面向 AI 的信任层（trust layer）——位于应用层防御（不可组合）与基础设施原语（缺乏工作流语义）之间的安全基础设施。

为解决 (1)，我们引入\textbf{认证工作流（authenticated workflows）}——一种协议级原语，在每次交互中强制执行\textit{意图}与\textit{完整性}。我们证明了四个攻击面是最小且完备的。协议级定位解决了异构性 (2) 并消除组合式缺口 (3)。面向 9 种主流框架的薄封装层解决了开发者负担 (4)，为企业 AI 信任层奠定基础。

\textbf{认证工作流}结合密码学、零信任（zero-trust）原则和运行时策略强制执行，确保每个操作真实、授权且可验证。我们提出一种新颖设计：使用嵌入各控制面（control surface）的分布式策略执行点（Policy Enforcement Point, PEP），独立验证密码学证明，以亚毫秒级开销运行，无需集中式基础设施。每个表面对操作独立验证后再执行，形成纵深防御——攻破一个表面不会攻破其他表面。

为指定\textit{意图}，我们引入 MAPL——一种新的 AI 原生策略语言，使用户能以动态、可扩展的方式表达智能体约束，即使智能体不断演进且调用上下文持续变化。MAPL 提供三项关键能力：\textit{可组合语法与继承}使策略能通过交集语义（intersection semantics）分层组合；\textit{极简规约}通过层级组合（hierarchical composition）与隐式主体（principal）解析将规模从 $O(M \text{ 用户} \times N \text{ 资源})$ 降至 $O(\log M + N)$；以及\textit{基于密码学证言（attestation）的动态状态}——证明操作已完成的签名声明，使策略得以表达如"仅在匿名化完成证言存在后才能导出数据"，并具备密码学强制力。这些进步克服了传统策略系统（OPA、AWS Cedar）在智能体环境中的局限。

本文聚焦跨操作边界的控制流保护——工具调用与数据检索——与 Rajagopalan 等人~\cite{rajagopalan2025prompts} 的研究互补，后者保护流经提示与上下文的数据。二者结合，为企业 AI 提供了形式化的信任层。通过一个支撑九种框架（第~\ref{sec:enforcement}节）的通用安全运行时，以薄适配器（200–500 LOC）实现，无需协议变更——证明了框架无关的安全能力。

形式上，我们证明了四个边界是完备且最小的，分布式强制执行是可靠的，策略组合保持安全属性。实验上，通过全面的攻击测试和生产 CVE 分析验证了这些属性。

当前智能体防御——护栏（guardrail）、提示过滤器、语义分析（模式匹配、ML 模型、启发式）——注重精确率但召回率低甚至为零。对抗性提示、编码技巧和新型模式通过语义盲区持续绕过它们，而误报则阻断合法操作。
与之对比，我们的密码学强制执行保证高精确率和高召回率：操作要么携带满足策略的有效签名，要么被直接拒绝。攻破系统需要击碎密码学原语——而非构造巧妙模式。这实现了完整的违规检测、零误报，以及基于计算困难性的确定性信任。


\noindent\textbf{贡献。}
我们提出认证工作流——一种协议级原语，在四个控制面（提示、工具、数据、上下文）上强制执行密码学验证，并证明其完备且最小（引理 6）。我们引入 MAPL，一种 AI 原生策略语言，通过密码学证言提供密码学工作流依赖，并通过层级组合将策略规约从 $O(M \times N)$ 降至 $O(\log M + N)$（定理 1-3）。我们展示了在九种异构框架（MCP、A2A、OpenAI、Claude、LangChain、CrewAI、AutoGen、LlamaIndex、Haystack）上的通用部署，通过薄适配器（200-500 LOC）实现，无需任何协议修改，分布式策略执行点嵌入每个边界提供零信任验证，开销为亚毫秒级。形式化证明确立了完备性与可靠性（引理 1-7）；实验验证展示了确定性保证（100\% 召回率，0\% 误报），覆盖 174 个测试用例，涵盖 10 项 OWASP Top 10 风险中的 9 项，并完全缓解两个生产 CVE（OpenAI Atlas、GitHub MCP）。
